% !TeX program = xelatex
% Text: CC BY-NC-SA 4.0. Companion code: Apache License 2.0.
\documentclass[UTF8,fontset=none,zihao=-4,a4paper]{ctexart}
\usepackage[margin=27mm]{geometry}
\usepackage{amsmath,amssymb,amsthm}
\usepackage{mathrsfs}
\DeclareFontFamily{U}{rsfs}{\skewchar\font127 }
\DeclareFontShape{U}{rsfs}{m}{n}{<5-7> rsfs5 <7-8> rsfs7 <8-> rsfs10}{}
\usepackage{hyperref}
\usepackage{xurl}
\setmainfont[Path=/mnt/c/Windows/Fonts/,BoldFont=timesbd.ttf,ItalicFont=timesi.ttf,BoldItalicFont=timesbi.ttf]{times.ttf}
\newfontfamily\arial[Path=/mnt/c/Windows/Fonts/]{arial.ttf}
\setCJKmainfont[Path=/mnt/c/Windows/Fonts/,BoldFont=simhei.ttf]{simsun.ttc}
\setCJKmonofont[Path=/mnt/c/Windows/Fonts/]{simsun.ttc}
\setCJKfamilyfont{hei}[Path=/mnt/c/Windows/Fonts/,AutoFakeBold=1.5]{simhei.ttf}
\newcommand\heiti{\CJKfamily{hei}}
\setlength{\parskip}{0pt}
\ctexset{
 section={format=\centering\heiti\bfseries\zihao{3},beforeskip=24bp,afterskip=18bp,aftername=\hspace{1em}},
 subsection={format=\raggedright\heiti\bfseries\zihao{4},beforeskip=24bp,afterskip=6bp,aftername=\hspace{1em}},
 subsubsection={format=\raggedright\heiti\bfseries\zihao{-4},beforeskip=12bp,afterskip=6bp,aftername=\hspace{1em}},
 paragraph={format=\raggedright\normalfont\zihao{-4},beforeskip=6bp,afterskip=6bp,aftername=\hspace{1em},runin=false}}
\setcounter{secnumdepth}{4}
\newtheorem{claim}{命题}
\newcommand{\OO}{\mathcal O}
\newcommand{\Q}{\mathbb Q}
\newcommand{\Z}{\mathbb Z}
\newcommand{\Vol}{\operatorname{Vol}}
\newcommand{\LogVol}{\operatorname{LogVol}}
\newcommand{\ord}{\operatorname{ord}}
\newcommand{\rad}{\operatorname{rad}}
\newcommand{\lana}[2]{\href{https://github.com/LANA-Project/iut4-sec1/blob/a7551d2f575b7e462cf9e82c587d9c6308e7b15c/#1}{\nolinkurl{#2}}}
\newcommand{\project}[2]{\href{https://github.com/NanachiKyubey/joshi-teich-abc/blob/47862524e9b23a21ca01cfeb112b342dffaac78a/#1}{\nolinkurl{#2}}}
\hypersetup{hidelinks,bookmarksnumbered=false,pdftitle={Kirti Joshi并未证出推论3.12和abc猜想},pdfauthor={Edison Mai}}
\begin{document}
\fontsize{12bp}{20bp}\selectfont
\begin{center}
{\heiti\zihao{1} Kirti Joshi 并未证出\\推论 3.12 和 abc 猜想\par}
\vspace{12bp}
Edison Mai\par
通讯作者 \href{mailto:old-baking-kung@duck.com}{old-baking-kung@duck.com}\par
2026 年 11 月 22 日
\end{center}
\vspace{24bp}
{\centering\heiti\zihao{3}\linespread{1}\selectfont 摘要\par}
\vspace{18bp}
\indent 截至 2026 年 11 月 22 日，望月新一已有数个月未能提供 LANA Project 形式化工作所需的推论 3.12。一个自然的想法是借助 Kirti Joshi 的算术 Teichmuller 空间论来形式化 LANA Project 所需的体积上界，并证明与该上界相配的下界，从而完成 \(abc\) 猜想的证明。本文证明，Joshi 系列第三篇中对应 3.12 的推论 9.11.1.1 按原文为假。在满足其初始条件且有奇坏处的数据上，推论要求一个非正数大于等于一个严格正数。即使订正符号，前一定理所用的加权体积也会给同一个张量格赋予两个不同的值。附录建立了修订模型中的有限处体积下界，并指出它与 Szpiro 型估计之间仍缺少什么。
\par\noindent\textbf{关键词}\quad 算术 Teichmuller 空间，推论 3.12，Tate 参数，Haar 测度，\(abc\) 猜想

\newpage
\vspace*{24bp}
{\centering\arial\zihao{3}\linespread{1}\selectfont Kirti Joshi Has Not Proved Corollary 3.12 or the \(abc\) Conjecture\par}
\vspace{18bp}
{\centering Edison Mai\par}
\vspace{8bp}
{\centering\arial ABSTRACT\par}
\vspace{6bp}
As of 22 November 2026, Shinichi Mochizuki had been unable for several months to supply the Corollary 3.12 input required for the LANA Project's formalization. A natural approach is to use Kirti Joshi's theory of arithmetic Teichmuller spaces to formalize the volume upper bound required by the LANA Project and to establish a matching lower bound, with the aim of completing a proof of the \(abc\) conjecture. We prove that Corollary 9.11.1.1 of Part III, Joshi's counterpart to Corollary 3.12, is false as stated. For data that satisfy its initial conditions and have an odd bad place, the corollary requires a nonpositive number to be at least a strictly positive number. Even after a sign correction, the weighted volume in the preceding theorem assigns two different values to the same tensor lattice. The appendix establishes a finite-place volume lower bound in a revised model and identifies what is still needed for a Szpiro-type estimate.
\par\vfill\noindent\textbf{KEY WORDS}\quad Arithmetic Teichmuller spaces, Corollary 3.12, Tate parameter, Haar measure, \(abc\) conjecture

\newpage
\section{绪论}
\subsection{IUT III 中的体积比较}
望月新一从数域上的单穿孔椭圆曲线、辅助素数 \(\ell\)、挠点覆盖和带标记的处出发，构造 Hodge theaters 及连接它们的 \(\Theta\)-link。IUT II 推论 4.10 给出 log-theta-lattice 的输入。IUT III 命题 3.9 固定对数体积的归一化，定理 3.11 给出多辐表示及其不定性。推论 3.12 取 \(\Theta\) 可能像的并集及其全纯包络，在论文规定的归一化和标签平均下把所得带符号对数体积记作 \(T\)。固定的 \(q\)-pilot 对数体积记作 \(x=-Q\)，其中 \(Q>0\)。目标为 \(x\le T\)。IUT IV 再用体积上界控制高度，因此上界和下界必须适用于同一个对象。\cite{mi,mii,miii,miv}

III 第 182--184 页先沿横向 link 识别抽象 pilot，再沿另一条路径计算 \(\Theta\) 输出包络的行列式和归一化全局度。第 (xi-e) 步指出，固定 \(q\)-pilot 的体积为 \(x\) 后，算法仍然可执行。第 (xi-f) 步由此写下 \(x\in(-\infty,T]\)。可执行性只给出输出的存在。要推出成员关系，还须证明这个输出的上端点至少为 \(x\)。把两条路径的数值放入同一个实数轴，或在 prime-strip 层次闭合同构链，仍需一条引理把固定的 \(q\)-pilot 度与另一条路径实际量到的包络相联系。\cite[第 181--185 页]{miii}

注 3.12.2(ii) 自己给出的实数模型可检验这一推断。取 \(h=1\)、\(\varepsilon=1/2\)，则输入和输出分别为
\[
x=-h=-1,\qquad (-\infty,-2h+\varepsilon]=(-\infty,-3/2].
\]
固定输入时可以执行 \(x\mapsto(-\infty,2x+\varepsilon]\)，但 \(x\) 不属于输出。若额外要求有同一个实数同时等于 \(-h\) 并属于输出，这个要求恰好等价于 \(h\le\varepsilon\)。该模型检验的是论证中的逻辑推断。满足 IUT 全部初始条件的算术反例仍未由此给出。一条充分的补证是，在实际 \(\Theta\) 包络内对每个 \(\delta>0\) 找到归一化全局对数体积至少为 \(x-\delta\) 的可测区域，然后用单调性令 \(\delta\) 趋于零。\cite[注 3.12.2(ii)，第 188--191 页]{miii}

\subsection{一个可直接复核的局部等式}
IUT III 注 3.9.5 的例子取 \(P=\Z_p^2\cup(\{p^{-1}\}\times\Z_p)\)，其 Haar 测度为 \(1\)，全纯包络为 \(p^{-1}\Z_p\times\Z_p\)。位于全纯包络内、测度同为 \(1\) 的分量理想只有 \(\Z_p^2\) 和 \(p^{-1}\Z_p\times p\Z_p\)。这是因为指数 \((a,b)\) 必须满足 \(a\ge-1,\ b\ge0,\ a+b=0\)。它们的并的测度为 \(2-p^{-1}\)，全纯包络的测度为 \(p\)。原文用底层 \(\Z_p\)-模自同构声称两者相等，但这样的自同构可把分量理想送到斜格。此处可直接订正，主推论的体积比较仍须独立证明。\cite[注 3.9.5，第 127--129 页]{miii}

\subsection{LANA Project 所需的数学输入}
LANA 的公开仓库形式化 IUT IV 第 1 节中可独立处理的部分，将 IUT I--III 的输入保留为显式接口或证书。本文还参照了一份本地 Lean 参考实现，其中 \nolinkurl{classicalABC_of_variant_genuine'} 由全称的具体 3.12 变体推出经典 \(abc\)。使用这个条件定理需要证明其前提，并核对它量到的具体区域与论文实际区域之间的比较。下文考察 Joshi III 和 IV 能否供给这个输入。\cite{lana,localref,jiii,jiv}

\section{Joshi 推论 9.11.1.1 的反例}
\subsection{原文不等式的符号}
令 \(m=(\ell-1)/2\)，令 \(S\) 为模域各处选出的代表中具有奇剩余特征和分裂乘法约化的处。写 \(a_w=|q_w^{1/(2\ell)}|_{L'_w}\)，其中 \(q_w\) 为局部 Tate 参数。Tate 一致化给出 \(0<a_w<1\)。记 \(V=\Vol(\widetilde\Theta_{\rm Mochizuki}^{\mathscr I})\)。定理 9.11.1（Form 2）声称 \(V\ge\prod_{w\in S}a_w^m\)。如果体积定义有效，Form 2 取对数得到
\begin{equation}
 m^{-1}\log V\ge \sum_{w\in S}\log a_w. \label{eq:form2}
\end{equation}
推论 9.11.1.1（源文件标签 \texttt{cor:cor312}）却印为
\begin{equation}
 -m^{-1}\big|\LogVol(\widetilde\Theta_{\rm Mochizuki}^{\mathscr I})\big|
 \ \ge\ -\sum_{w\in S}\log a_w. \label{eq:printed}
\end{equation}
若 \(S\ne\varnothing\)，式 \eqref{eq:printed} 的右端严格为正，左端在体积有实对数时至多为零。因此无论体积取什么值，都得不到原文结论。在局部对数非正且可求和的约定下，可把式 \eqref{eq:form2} 的右端写成 \(-\sum_w|\log a_w|\)。Joshi II 的局部版本和 IV 调用 III 时均使用两个带负号的绝对值之间的比较。\cite[定理 9.11.1 至推论 9.11.1.1]{jiii}\cite{jii,jiv}

\subsection{满足初始条件的曲线}
取 \(k=\Q(\alpha)\)、\(\alpha^3=2\)、\(J=\alpha/5\)。\(X^3-2\) 由 Eisenstein 判别法不可约，判别式为 \(-108\)。它模 \(5\) 分解为 \((X-3)(X^2+3X+4)\)，其中二次因子的判别式为 \(3\bmod 5\)，故不可约。于是 \(5\) 在 \(k\) 中有两个不分歧的处 \(v_1,v_2\)，局部次数分别为 \(1,2\)，且 \(\ord_{v_i}(J)=-1\)。在 \(k\) 上取
\[
 C\colon y^2+xy=x^3-\frac{36}{J-1728}x-\frac1{J-1728}.
\]
有 \(c_4=J/(J-1728)\)、\(\Delta=J^2/(J-1728)^3\)，故 \(j(C)=J\)，模域为 \(k\)。两处的方程均整，\(c_4\) 为单位，\(\ord_{v_i}(\Delta)=1\)。约化曲线 \(y^2+xy=x^3\) 的节点有两条有理切线 \(y=0,-x\)，故两处均为分裂乘法约化，这个性质在有限扩张后保持。

取 \(L=k(i,C[6])\)，它在 \(k\) 上有限 Galois，且全虚并含全部 \(6\) 阶挠点。单穿孔曲线 \(C-\{O\}\) 属于文中所用的严格 Belyi 型情形。由于 \(J\) 在 \(5\) 处有极点，而复乘 \(j\) 不变量为代数整数，\(C\) 无复乘。Serre 开像定理给出所有充分大素数 \(\ell\) 上包含 \(\mathrm{SL}_2(\mathbb F_\ell)\) 的模 \(\ell\) 像。取这样的 \(\ell>5\) 且 \(\ell\nmid[L:k]\)，令 \(L'=L(C[\ell])\)，并为每一模域处选一个 \(L'\) 上方的代表。\(v_1,v_2\) 的代表 \(w_1,w_2\) 均属于 \(S\)，故式 \eqref{eq:printed} 在该初始数据上失败。\cite[III 的初始数据段]{jiii}\cite{jii,serre}

\subsection{加权体积没有定义在张量子集上}
对有限扩张 \(E_i/\Q_p\)，令 \(\mu_i(\OO_{E_i})=1\)，\(d_i=[E_i:\Q_p]\)。III 规定张量格的体积为
\begin{equation}
 \Vol^\Gamma\left(\bigotimes_i V_i\right)=\prod_i\mu_i(V_i)^{\gamma_i},
 \qquad \gamma_i=[L'_{w_i}:k_{v_i}]^{-1} \quad(E_i=L'_{w_i}).
 \label{eq:weights}
\end{equation}
这里张量格是 \(\Q_p\)-张量积的子集，整格的张量在 \(\Z_p\) 上取。按处展开张量容器，包含 \(E_1\otimes_{\Q_p}E_2\) 这样的混合分量。写 \(M=\OO_{E_1}\otimes_{\Z_p}\OO_{E_2}\)。张量积的平衡关系给出同一个子集的两种表示
\[
 (p\OO_{E_1})\otimes_{\Z_p}\OO_{E_2}=pM
 =\OO_{E_1}\otimes_{\Z_p}(p\OO_{E_2}).
\]
因为 \([\OO_{E_i}:p\OO_{E_i}]=p^{d_i}\)，两种表示分别给出 \(p^{-d_1\gamma_1}\) 和 \(p^{-d_2\gamma_2}\)。由塔式公式 \(d_i\gamma_i=[k_{v_i}:\Q_p]\)。代入上一节的处，得到
\[
 \Vol^\Gamma(5M)=5^{-1},\qquad \Vol^\Gamma(5M)=5^{-2}.
\]
在更多槽的张量中，其余槽取整格，仍得到同一矛盾。式 \eqref{eq:weights} 可以给带标签的因子格元组赋值。要把它用于子集的包络，就必须同意一个子集的所有张量表示。改正式 \eqref{eq:printed} 的符号也不会改变上述计算。\cite[III 的带权体积定义及 Form 2]{jiii}

这里的 \(\mu_i\) 是 Haar 测度。若 \(|p|_{E_i}=p^{-1}\)，则 \(\mu_i(a\OO_{E_i})=|N_{E_i/\Q_p}(a)|_p=|a|_{E_i}^{d_i}\)。用绝对值重写计算时必须带上局部次数。上面的反例直接使用格指数，独立于绝对值的记号。

\section{从局部构造到 Form 2 的证明链}
\subsection{赋值尺度与周期环的上确界}
Joshi I 从 perfectoid 域的 untilt 与几何基点出发构造算术全纯结构。II 使用 Fargues--Fontaine 曲线上的点 \([a^{j^2}]-p\)，其中 \(0<|a|<1\)。在各自从 tilt 诱导的赋值下有 \(v_j(p)=j^2v_1(p)\)，故以第 \(m\) 个分量为基准时有 \(v_j(p)=j^2v_m(p)/m^2\)。若同时要求同一个赋值域等距嵌入所有分量，则各分量中 \(p\) 的数值相同。II.5 的全局归一化恢复积公式，应用局部平方缩放前须说明使用的是哪一套尺度。III 的 Form 1 证明把 \(j^2/m^2\) 用在第一个分量上。代 \(j=1\) 给出 \(a=a^{1/m^2}\)，在 \(0<a<1\)、\(m>1\) 时不成立。\cite{ji,jii,jhalf,jiii}

如果各提升确实进入指定的周期环集合，并使用以第 \(m\) 项为标准的诱导尺度，设根在该处的范数为 \(a\in(0,1)\)，则一组可用的实数计算为
\[
 \prod_{j=1}^m a^{j^2/m^2}
 =a^{(m+1)(2m+1)/(6m)}\ge a^m.
\]
有限个坏处相乘并在包含这些见证的集合上取上确界，就给出周期环版本的估计。上同调容器的体积需要另行比较。\cite[III 的周期环估计和 Form 1]{jiii}

\subsection{原级数的 \texorpdfstring{\(\theta\)}{theta} 值}
Joshi II 所印级数约去前因子后为
\[
 \theta_M(u)=\sum_{n\in\Z}q^{n(n+1)/2}u^{2n+1}.
\]
换指标 \(n\mapsto-n-1\) 给出 \(\theta_M(u^{-1})=\theta_M(u)\)。在奇剩余特征下，\(\theta_M(1)\) 的常数贡献为 \(2\)，所以不为零。取 \(p\ne\ell\) 及本原 \(\ell\) 次单位根 \(\zeta\)。在 \(u=q^{j/(2\ell)}\zeta\) 处，项的 \(q\) 指数为 \(b_n=n(n+1)/2+(2n+1)j/(2\ell)\)。由于 \(b_{n+1}-b_n=n+1+j/\ell\)，唯一最小项在 \(n=-1\)，指数 \(-j/(2\ell)\)。分母 \(\theta_M(\zeta)\) 的首项和 \(\zeta+\zeta^{-1}\) 模极大理想非零，因为 \(\bar\zeta\) 的阶是奇素数。因此这个级数给出比值范数 \(|q|^{-j/(2\ell)}\)，原文采用的是 \(|q|^{-j^2/(2\ell)}\)。取 \(j=2\) 即有差别。即使直接用 Tate 根代替该级数，其范数 \(|q|^{1/(2\ell)}\) 也与同一绝对值下的原 Tate 参数范数不同。\cite[II 的 theta 函数与局部取值]{jii}\cite{jiii}

\subsection{上同调、张量关系与包络}
给定非零整元 \(r=q^{1/(2\ell)}\)，置 \(p^*=p\)（奇 \(p\)）或 \(p^*=4\)（\(p=2\)）。用主单位 \(1+p^*r\) 的相容根塔可构造单位 Kummer 类。通常 Bloch--Kato 对数在这里对应 \(\log_p(1+p^*r)\)。除以 \(p^*\) 得到一个新的壳坐标。III 同时把这个缩放后的坐标称为通常指数的逆。以 \(u=1+p\in\Q_p\) 为例，通常对数的阶为 \(1\)，归一化坐标的阶为 \(0\)。区分这两个映射后，可保留附录中使用的等距计算。\cite[III 的 Kummer 类和 log-shell]{jiii}

将不同 arithmeticoid 的类放入共同容器后，纯张量映射 \((u,v)\mapsto u\otimes v\) 可以定义集合的像。它在直积上不保持加法，因为 \((1,0)\) 和 \((0,1)\) 的像均为零，而和的像 \(1\otimes1\) 非零。交叉范数也必须遵守 \(p\otimes1=1\otimes p\)。若 \(\|u\otimes v\|=|u|_1|v|_2\)，则这条等式要求 \(|p|_1=|p|_2\)，与在同一张量中保留平方缩放矛盾。相关的上同调同构与全局 Frobenius 移位可以组织见证，仍须在统一标量规范下重算它们的大小。\cite{jhalf,jiii,jiv,janab,jalg}

全纯包络与普通非阿基米德凸闭包也会不同。取奇 \(p\)，\(E=\Q_p(\pi)\)、\(\pi^2=p\)。在 \(E\otimes_{\Q_p}E\simeq E\times E\)、\(a\otimes b\mapsto(ab,a\bar b)\) 下，整张量格的像 \(P\) 已经是紧的 \(\Z_p\)-格，因而是闭凸的。它的两坐标模 \(\pi\) 相同，因为共轭在剩余域上恒等，所以 \((1,0)\notin P\)。但两坐标投影均为 \(\OO_E\)，其最小分量整理想包络为 \(\OO_E\times\OO_E\)。Form 2 的证明从含一个标准见证过渡到含整个张量格，正需要这样的额外闭包操作。若实际集合 \(U\) 含整点 \(t\)，其逐分量理想包络确实含有 \(t\OO_T\)，其中 \(\OO_T\) 是张量代数的整闭包。附录用这项包含重新计算体积。\cite[III 的 holomorphic hull、体积引理和 Form 2]{jiii}

\subsection{第四篇对第三篇的依赖}
Joshi IV 在从体积比较到高度估计时调用 III 的下界。原文推论在上述算术实例上为假，而 Form 2 所需的子集体积在同一实例上无法满足其定义。IV 中改用带负号的绝对值并移位标准点，仍须另证良定义的体积及对实际包络的下界。因此所审阅的 ATS 证明链没有建立其宣称的 \(abc\) 结论。\cite{jiii,jiv}

\appendix
\section{从 ATS 回收有限处体积下界}
\subsection{在整闭包上量体积}
设 \(T\) 是有限多个有限 \(\Q_p\)-扩张在 \(\Q_p\) 上的张量积。它是有限 étale 代数，故 \(T=\prod_\alpha F_\alpha\)，且 \(\OO_T=\prod_\alpha\OO_{F_\alpha}\)。以 \(\mu_T(\OO_T)=1\) 归一化加法 Haar 测度，记 \(\nu_T(U)=(\dim_{\Q_p}T)^{-1}\log\mu_T(U)\)。对于有界整集 \(U\)，若每个分量至少含一个非零投影，则其最小分量理想包络是
\[
 H(U)=\prod_\alpha\pi_\alpha^{n_\alpha}\OO_{F_\alpha},
 \qquad n_\alpha=\min_{u\in U}\ord_\alpha(u_\alpha).
\]
离散性保证每个最小值能取到。对单点 \(t\) 有 \(H(\{t\})=t\OO_T\)，其体积由乘法行列式给出。若 \(T=\bigotimes_i E_i\) 且 \(t=\bigotimes_i t_i\) 的各 \(t_i\ne0\)，张量行列式公式给出
\begin{equation}
 \nu_T(t\OO_T)=\sum_i\log|t_i|_{\rm st},
 \qquad |p|_{\rm st}=p^{-1}. \label{eq:haarnorm}
\end{equation}
具体地，第 \(i\) 槽的乘法行列式为 \(N_{E_i/\Q_p}(t_i)\)，在整个张量中的幂次是 \(\dim T/[E_i:\Q_p]\)。除以 \(\dim T\) 得式 \eqref{eq:haarnorm}。因为左端是同一张量代数上的行列式，此计算与张量表示的选择相容。

\subsection{单个有效坐标的 Kummer 见证}
设 \(K/F\) 为含所需局部 Tate 根的有限数域扩张，\(B_F\) 为选定的有限坏处集。在 \(s\mid w\) 且 \(w\in B_F\) 时取 \(r_s^{2\ell}=q_w\)，其余处取 \(r_s=1\)。定义
\[
 \tau_p(r)=p^{*-1}\log_p(1+p^*r).
\]
对非零整元 \(r\)，有 \(v_p(p^*r)>1/(p-1)\)。由于 \(v_p(n)\le(n-1)/(p-1)\)，对数级数的第一项严格支配其余各项。所以 \(|\tau_p(r)|_{\rm st}=|r|_{\rm st}\)，尤其 \(\tau_p(1)\) 是单位。

\begin{claim}
在第 \(j\) 个标签组的 \(j+1\) 槽张量中，仅在第 \(j\) 槽放入 \(\tau_p(r_s)\)，其余槽放入 \(\tau_p(1)\)。将所得见证点取逐分量主理想包络，并对各标签组做次数归一化，有限处的平均对数体积为 \(-Q\)，其中
\[
 Q=\sum_p Q_p,\qquad
 Q_p=\frac{\log p}{2\ell[F:\Q]}
      \sum_{\substack{w\mid p\\w\in B_F}} f_w\ord_w(q_w).
\]
这里 \(f_w\) 是 \(F_w/\Q_p\) 的剩余次数。包含这些理想的有限处区域也有相应的下界，只要其全局体积可定义。
\end{claim}
\begin{proof}
写 \(A_p=K\otimes_{\Q}\Q_p=\prod_{s\mid p}K_s\)，\(\lambda_s=[K_s:\Q_p]/[K:\Q]\)。按处元组 \(c=(s_0,\ldots,s_j)\) 展开 \(A_p^{\otimes(j+1)}\)。在 \(c\) 分量使用式 \eqref{eq:haarnorm}，得到归一化对数体积 \(\log|r_{s_j}|_{\rm st}\)。该分量在整个张量代数中的维数占比是 \(\prod_i\lambda_{s_i}\)。由于 \(\sum_{s\mid p}\lambda_s=1\)，对所有填充槽求和后只剩
\[
 h_{j,p}=\sum_{s\mid p}\lambda_s\log|r_s|_{\rm st}.
\]
对 \(s\mid w\in B_F\)，令 \(e_w\) 为 \(F_w/\Q_p\) 的分歧指数。此时
\[
 -\log|r_s|_{\rm st}=\frac{\ord_w(q_w)}{2\ell e_w}\log p,\qquad
 \sum_{s\mid w}\lambda_s=\frac{e_wf_w}{[F:\Q]}.
\]
后一等式来自 \(\sum_{s\mid w}[K_s:F_w]=[K:F]\)。故每个标签组均有 \(h_{j,p}=-Q_p\)，对 \(p\) 求和并取标签平均即得结论。其余处的贡献为零，包含关系的结论由 Haar 测度单调性得到。
\end{proof}

\subsection{与参考区域的差别及其后果}
具体参考区域在有效槽使用 \(r_s^{j^2}\)。允许的局部域自同构保持赋值，因而保持该主理想。对 \(j>1\) 和 \(0<|r_s|_{\rm st}<1\)，有 \(r_s^{j^2}\OO\subsetneq r_s\OO\)，上述 Kummer 见证在相应分量落在较大的理想中。次数加权后的参考有限处体积为 \(-j^2Q_p\)。在具体参考无穷处体积为 \((j+1)\log\pi\) 的模型中，平均全局体积是
\[
 R_{\rm ref}=-A_mQ+B_m,\qquad
 A_m=\frac{(m+1)(2m+1)}6,\qquad B_m=\frac{m+3}2\log\pi .
\]
所以 \(-Q\le R_{\rm ref}\) 等价于 \((A_m-1)Q\le B_m\)。对 \(\ell=2m+1\ge7\)，它要求 \(2\ell Q\le(126/11)\log\pi\)。例如有理 Legendre 参数 \(x_N=5^N/(2\cdot5^N+1)\) 的 \(j\) 不变量在 \(5\) 处有 \(2N\) 阶极点。当全部 \(5\) 上方的处进入支持时，\(2\ell Q\ge2N\log5\)，故 \(N\ge6\) 时超过阈值。对应本地存在定理还显式接收曲线数据、有限性和远阿贝尔几何存在性输入。此计算定位的是具体参考区域中的障碍。\cite{localref,adapt}

在 ATS 的 Ansatz 中，可以尝试对每个标签组单独选择有效坐标所在的标准点，以产生一次幂见证。仍需把见证送入上界实际所量的同一包络，并控制无穷处和各处的权重。在下降到较小的局部域时，亦可对整个整理想取迹。对局部扩张 \(E/k\)，若 \(\mathfrak D_{E/k}=\mathfrak p_E^d\)，则
\[
 {\rm Tr}_{E/k}(\mathfrak p_E^n)
   =\mathfrak p_k^{\lfloor(n+d)/e(E/k)\rfloor}.
\]
这是因为 \({\rm Tr}(\mathfrak p_E^n)\subseteq \mathfrak p_k^a\) 等价于
\(\mathfrak p_E^{n-ea}\subseteq\mathfrak D_{E/k}^{-1}\)。若 \(b\OO_E=\mathfrak p_E^d\)，用 \({\rm Tr}(b^{-1}\mathfrak p_E^n)\) 并向上取整到 \(\mathfrak p_k^{\lceil n/e\rceil}\)，相对于 \(n/e\) 的额外指数小于 \(1\)。在有限坏处按全局次数加权，这种取整损失至多为坏处根基级的对数和。这个计算控制的是整层理想。一个单独见证的迹可因抵消而更小，现有的逐元组上界也没有因此自动适用于新理想。\cite{riemcode,adapt}

\subsection{从下界到 Szpiro 和 \texorpdfstring{\(abc\)}{abc} 的条件}
新下界仍带有每个 Tate 参数的阶数。固定奇素数 \(p\)，对 \(a=p^N,\ b=1,\ c=p^N+1\)，Frey 曲线 \(y^2=x(x-a)(x+b)\) 在 \(p\) 处的最小判别式阶和 Tate 参数阶均为 \(2N\)。而 \(p\) 对 \(\log\rad(abc)\) 的贡献一直是 \(\log p\)。\(p^N+1\) 的其他素因子也进入全局根基。要证明某一族曲线上的 Szpiro 型界，还须证明把这些重数与导子或根基相联系的估计。

若能在同一个实际包络上取得
\[
 -Q-\delta\le R\le-AQ+U,\qquad A>1,\quad \delta\ge0,
\]
便有 \((A-1)Q\le U+\delta\)。在适用族上把右端统一换算成根基和高度，才有可能得到该族的 Szpiro 或 \(abc\) 估计。若对全部互素正整数三元组 \(a+b=c\) 推出 \(\log c\le B\log\rad(abc)+C\)，其中 \(B>1,C\) 固定，便得到固定指数的弱 \(abc\) 估计。完整 \(abc\) 需要对每个 \(\varepsilon>0\) 取得 \(B=1+\varepsilon\)，并让 \(C_\varepsilon\) 对全部三元组统一。本文证明的是有限处下界。Szpiro 型局部特例、固定指数和完整 \(abc\) 都还需要上述比较及相应的量词控制。

若上界的 \(A\) 取参考模型的平方权重平均 \(A_m=\ell(\ell+1)/12\)，记 \(L=2\ell Q\)，则消去 \(R\) 给出
\[
 \frac L6\le \frac{4\ell}{(\ell-3)(\ell+4)}(U+\delta).
\]
对有理 Frey 曲线，若所选坏处都在 \(2abc\) 上方，整层理想下降带来的误差至多为 \(\log\rad(abc)+\log 2\)。由于此项在上式中的系数为 \(O(1/\ell)\)，若能为所需三元组统一选取足够大的适用辅助素数，这个误差本身允许指数逼近 \(1\)。决定指数的其余工作仍包括实际几何包络的上下界以及 \(U\) 的算术控制。\par

\section{充分利用并扩张Dirichlet GRH的Theta提升技术以进一步挽救ATS}

2026 年 11 月 21 日，Lean 社区博客宣布已核验 Levent Alpöge 的经典黎曼猜想证明。Riemann 工程的 Lean 源码还给出了 Dirichlet \(L\) 函数在临界带内的零点定理，即零点的实部为 \(1/2\)。\cite{riemblog,riemcode}

该工程的迹对偶与不同理想公式可用于附录 A 的理想下降。显式素数计数可用于控制辅助素数 \(\ell\) 的选取误差，这两项应用均可独立于零点定理。这里进一步考察该工程采用的自守 Theta 提升及 Fourier 技术能否用来估计 ATS 的局部体积缺额。要将这一估计用于 Szpiro 或 \(abc\)，仍须在实际包络中实现相应的 Kummer 见证，并在同一包络上建立全局上界。\cite{riemcode,riemtheta}

\subsection{从相邻 Theta 值取得局部见证}
固定奇坏处 \(p\ne\ell\)。在含有 \(r=q^{1/(2\ell)}\) 和本原 \(\ell\) 次单位根 \(\zeta\) 的同一个局部域 \(E\) 中，按 Joshi II 印出的级数定义
\[
 \xi_0=1,\qquad \xi_j=\frac{\theta_M(\zeta)}{\theta_M(r^j\zeta)},\qquad
 u_j=\xi_j/\xi_{j-1}.
\]
由正文的首项计算，\(v(\xi_j)=jv(r)\)，因此 \(v(u_j)=v(r)\)。附录 A 的对数等距性给出 \(\tau_p(u_j)\OO_E=r\OO_E\)。这为一次幂理想提供了来自 Theta 值的生成元。用于 ATS 时，须证明 \(1+p^*u_j\) 的 Kummer 类经 III 允许的汇集进入指定标签。两个 Kummer 类相减对应的是其主单位之商。跨 untilt 作商还需要相容的共同域和标记。\cite{jii,jiii}

\subsection{一个可计算的局部生成例子}
取 \(p=29,\ell=7,q=29,\pi^{14}=29\)，并取 \(\zeta\in\Q_{29}\) 为本原七次单位根。在 \(E=\Q_{29}(\pi)\) 中令 \(W=\pi\OO_E\)、\(L_0=\pi^4\OO_E\)，对应标签 \(j=2\)。\(W/L_0\) 有三个 \(\mathbb F_{29}\) 方向。十四个根分支 \(\omega\pi\) 的缩放对数模 \(L_0\) 都在线 \(\mathbb F_{29}\pi\) 上，因为 \(\mu_{14}\subset\Q_{29}\)。加入 \(L_0\) 后，生成的格在 \(W\) 中仍有指数 \(29^2\)。

直接展开级数，保留模 \(\pi^4\) 所需的项，得到
\[
 \tau_{29}(\pi)\equiv\pi,\qquad
 \tau_{29}(\xi_2)\equiv(1+\zeta^2)\pi^2,\qquad
 \tau_{29}(u_2)-\tau_{29}(\pi)\equiv\zeta^2\pi^3.
\]
最后一式来自 \(\theta_M(\pi\zeta)/\theta_M(\pi^2\zeta)\equiv\pi+\zeta^2\pi^3\)。缩放对数的高次项模 \(\pi^4\) 消失，且 \(1+\zeta^2\) 是单位。三个指定点连同 \(L_0\) 因而生成整个 \(W\)。应用此等式需要在同一实际标签中实现三个点，并控制其余张量槽。若改取 \(r_t=29^t\pi\)（\(t\ge1\)），则 \(r_t^4\OO_E\subset29r_t\OO_E\)，而 \(r_t\OO_E/29r_t\OO_E\) 有十四个方向。任取三个点也不足以通过同样的生成论证覆盖这族数据。

\subsection{用 Fourier 对偶衡量体积缺额}
在完整的 \((j,p)\) 张量代数中记维数为 \(d\)，沿用附录 A 的体积归一化。一次幂格 \(W\) 与平方幂格 \(L_0\) 分别满足 \(\nu(W)=-Q_p\)、\(\nu(L_0)=-j^2Q_p\)。若实际局部区域 \(C\) 是包含 \(L_0\) 的可测加法子群，令 \(H=(C\cap W)/L_0\subset W/L_0\)。由有限群对偶，\(H\) 的零化子 \(H^\perp\) 有 \([W:C\cap W]\) 个字符，故
\[
 \nu(C)\ge -Q_p-\frac{\log|H^\perp|}{d}.
\]
对于一般可测区域，只要占比 \(\rho=\mu(C\cap W)/\mu(W)>0\)，也有 \(\nu(C)\ge -Q_p+d^{-1}\log\rho\)。因此可以按全局归一化汇总各处的缺额。

Riemann 工程的 Heisenberg 平移展开、Weil 表示的 Fourier 公式和迹对偶提供了处理指定局部陪集及其字符的工具。现有 Rallis 公式在给定局部条件后，存在性地选择测试多项式、测度、基本域和常数，计算提升的 Petersson 平方范数。要用于上式，须从 III 实际允许的 Kummer 像构造指定测试数据，并证明其 Fourier 系数控制 \(H^\perp\) 或 \(\rho\)。若用极化得到混合内积，还须在相同归一化下控制两个范数相减后的误差。Dirichlet GRH 的零点位置结论本身没有提供这些比较。\cite{riemtheta}

\subsection{线性权重足以保留首阶系数}
若只能实现单个 \(\xi_j\)，则 \(\tau_p(\xi_j)\OO_E=r^j\OO_E\)。在单有效槽模型按次数加权得到 \(-jQ_p\)，前提是各坏处都有相应见证。记 \(b_m=(\ell+1)/4\) 为线性平均，\(A_m=\ell(\ell+1)/12\) 为平方平均。若在同一实际全局包络上取得
\[
 -b_mQ-\delta\le R\le-A_mQ+U,
 \qquad
 \frac{L}{6}\le \frac{4\ell}{(\ell-3)(\ell+1)}(U+\delta),
 \quad L=2\ell Q,\quad\ell>3,
\]
则第二式由第一式消去 \(R\) 得出。相对附录 A 的一次幂方案，系数只增加因子 \((\ell+4)/(\ell+1)\)，仍有 \(4/\ell\) 的首阶项。

{\setlength{\emergencystretch}{2em}在 \(p=337,\ell=7,j=2\) 的指定处元组上，线性格的对数体积为 \(-(6N+30/7)\log337\)，已有逐元组上界为 \(-(12N+5)\log337\)。差额 \((6N+5/7)\log337\) 随 \(N\) 无界。要由这一路线推出 Szpiro 型估计，仍须证明实际包络中见证的覆盖程度，以及对同一包络成立的全局上界或跨处补偿。\par}

\section*{人工智能使用声明与许可}
Edison Mai 主导研究并任通讯作者。OpenAI ChatGPT（Astra 6.2）协助阅读原文和代码，整理数值反例及修订模型的自然语言推导，并起草、排版和检查本文。这篇文章和自然语言材料采用 \href{https://creativecommons.org/licenses/by-nc-sa/4.0/}{CC BY-NC-SA 4.0} 许可。配套代码采用 \href{https://www.apache.org/licenses/LICENSE-2.0}{Apache License 2.0}。

\clearpage
\begingroup
\renewcommand{\heiti}{\rmfamily}
\ctexset{section/name={,},section/number=\arabic{section},
 appendix/name={Appendix\space,},appendix/number=\Alph{section}}
\markboth{}{}
\renewcommand{\proofname}{Proof}
\setcounter{section}{0}
\setcounter{claim}{0}
\renewcommand{\thesection}{\arabic{section}}
\renewcommand{\theclaim}{\arabic{claim}}
\renewcommand{\theHsection}{english.\arabic{section}}
\renewcommand{\theHclaim}{english.\arabic{claim}}
\setcounter{equation}{0}
\renewcommand{\theHequation}{english.\arabic{equation}}
\begin{center}
{\arial\zihao{3} Kirti Joshi Has Not Proved Corollary 3.12 or the \(abc\) Conjecture\par}
\vspace{18bp}
Edison Mai
\end{center}

\section{Introduction}
\subsection{The volume comparison in IUT III}
Mochizuki starts with a once-punctured elliptic curve over a number field, an auxiliary prime \(\ell\), torsion covers, and marked places. He constructs Hodge theaters and the \(\Theta\)-link between them. Corollary 4.10 of IUT II supplies data for the log-theta-lattice. Proposition 3.9 of IUT III normalizes log-volume, and Theorem 3.11 gives the multiradial representation and its indeterminacies. In Corollary 3.12, write \(T\) for the signed log-volume of the holomorphic hull of the union of possible \(\Theta\)-images, with the prescribed normalization and label average. Write \(x=-Q\), \(Q>0\), for the fixed \(q\)-pilot log-volume. The desired comparison is \(x\le T\). IUT IV combines it with an upper bound to control heights. The two bounds must apply to the same object.\cite{mi,mii,miii,miv}

Pages 182--184 of IUT III identify the abstract pilots along the horizontal link. A separate construction computes the determinant and normalized global degree of the output hull. Step (xi-e) says that the construction remains executable with the input volume fixed at \(x\). Step (xi-f) then asserts \(x\in(-\infty,T]\). Executability gives an output. To establish membership, one also needs a lower bound of \(x\) for its upper endpoint. A common real coordinate and a closed chain of prime-strip isomorphisms still require a comparison between the fixed degree and the hull actually produced.\cite[pp. 181--185]{miii}

Remark 3.12.2(ii) gives a real-number model. Take \(h=1\) and \(\varepsilon=1/2\). The input is \(x=-h=-1\), whereas the output interval is \((-\infty,-2h+\varepsilon]=(-\infty,-3/2]\). The map \(x\mapsto(-\infty,2x+\varepsilon]\) is executable for this fixed input, and \(x\) lies outside its output. Requiring one real number to have both properties is equivalent to requiring \(h\le\varepsilon\). This tests the inference in the proof. An arithmetic example satisfying all IUT initial conditions would require additional work. One sufficient way to establish the missing comparison is to construct, for every \(\delta>0\), a measurable region inside the actual hull with normalized global log-volume at least \(x-\delta\). Monotonicity then gives \(x\le T\).\cite[Remark 3.12.2(ii), pp. 188--191]{miii}

\subsection{A local equality that can be checked directly}
In Remark 3.9.5 of IUT III, the set \(P=\Z_p^2\cup(\{p^{-1}\}\times\Z_p)\) has Haar measure \(1\) and hull \(p^{-1}\Z_p\times\Z_p\). A product ideal \(p^a\Z_p\times p^b\Z_p\) inside the hull has measure \(1\) precisely when \(a\ge-1\), \(b\ge0\), and \(a+b=0\). The only possibilities are \(\Z_p^2\) and \(p^{-1}\Z_p\times p\Z_p\). Their union has measure \(2-p^{-1}\), while the hull has measure \(p\). The text uses a \(\Z_p\)-module automorphism to assert equality. Such an automorphism can take a product ideal to a skew lattice. The main comparison still requires a separate proof after this correction.\cite[Remark 3.9.5, pp. 127--129]{miii}

\subsection{The input required by the LANA Project}
The public LANA repository formalizes parts of IUT IV, Section 1, and keeps inputs from IUT I--III as explicit interfaces or certificates. In the local Lean reference implementation consulted here, \nolinkurl{classicalABC_of_variant_genuine'} derives classical \(abc\) from a universally quantified concrete variant of Corollary 3.12. Applying it requires a proof of that premise and a comparison between its concrete region and the region in the paper. We examine whether Joshi III and IV supply this input.\cite{lana,localref,jiii,jiv}

\section{A Counterexample to Joshi's Corollary 9.11.1.1}
\subsection{The signs in the printed inequality}
Put \(m=(\ell-1)/2\). Let \(S\) be the chosen representatives of odd residue characteristic with split multiplicative reduction. At \(w\in S\) write \(a_w=|q_w^{1/(2\ell)}|_{L'_w}\), where \(q_w\) is the local Tate parameter. Tate uniformization gives \(0<a_w<1\). Put \(V=\Vol(\widetilde\Theta_{\rm Mochizuki}^{\mathscr I})\). Theorem 9.11.1 (Form 2) asserts \(V\ge\prod_{w\in S}a_w^m\). If this volume is defined, taking logarithms gives
\begin{equation}
 m^{-1}\log V\ge\sum_{w\in S}\log a_w. \label{eq:en-form2}
\end{equation}
Corollary 9.11.1.1, labeled \texttt{cor:cor312} in the source, instead prints
\begin{equation}
 -m^{-1}\big|\LogVol(\widetilde\Theta_{\rm Mochizuki}^{\mathscr I})\big|
 \ge -\sum_{w\in S}\log a_w. \label{eq:en-printed}
\end{equation}
For nonempty \(S\), its right-hand side is strictly positive and its left-hand side is nonpositive for every real log-volume. Under the convention that the local log-volumes are nonpositive and summable, the right-hand side of \eqref{eq:en-form2} can be written as \(-\sum_w|\log a_w|\). The local variant in Part II and the invocation of Part III in Part IV both compare two negative absolute values.\cite{jii,jiii,jiv}

\subsection{A curve satisfying the initial hypotheses}
Take \(k=\Q(\alpha)\), \(\alpha^3=2\), and \(J=\alpha/5\). The polynomial \(X^3-2\) is Eisenstein at \(2\) and has discriminant \(-108\). Modulo \(5\) it factors as \((X-3)(X^2+3X+4)\). The quadratic factor is irreducible because its discriminant is \(3\bmod 5\). Thus \(5\) has two unramified places \(v_1,v_2\) in \(k\), of local degrees \(1,2\), and \(\ord_{v_i}(J)=-1\). Over \(k\) take
\[
 C\colon y^2+xy=x^3-\frac{36}{J-1728}x-\frac1{J-1728}.
\]
Its invariants are \(c_4=J/(J-1728)\) and \(\Delta=J^2/(J-1728)^3\), so \(j(C)=J\) and the field of moduli is \(k\). At both places the equation is integral, \(c_4\) is a unit, and \(\ord_{v_i}(\Delta)=1\). The reduction is \(y^2+xy=x^3\), with rational tangent lines \(y=0,-x\) at the node. The reduction is split multiplicative at both places and stays so after finite extension.

Let \(L=k(i,C[6])\). This is a totally imaginary finite Galois extension of \(k\) over which all \(6\)-torsion is rational. The once-punctured curve \(C-\{O\}\) is of the strict Belyi type used in Part III. The pole of \(J\) at \(5\) excludes complex multiplication because CM \(j\)-invariants are algebraic integers. Serre's open image theorem gives a mod-\(\ell\) image containing \(\mathrm{SL}_2(\mathbb F_\ell)\) for all sufficiently large primes \(\ell\). Choose such an \(\ell>5\) with \(\ell\nmid[L:k]\), let \(L'=L(C[\ell])\), and choose one place of \(L'\) above every place of \(k\). The representatives \(w_1,w_2\) of \(v_1,v_2\) both belong to \(S\). Thus \eqref{eq:en-printed} fails on admissible initial data.\cite{jii,jiii,serre}

\subsection{The weighted volume does not descend to tensor subsets}
For finite extensions \(E_i/\Q_p\), let \(\mu_i(\OO_{E_i})=1\) and \(d_i=[E_i:\Q_p]\). Part III prescribes the rule
\begin{equation}
 \Vol^\Gamma\left(\bigotimes_i V_i\right)=\prod_i\mu_i(V_i)^{\gamma_i},
 \qquad \gamma_i=[L'_{w_i}:k_{v_i}]^{-1},\quad E_i=L'_{w_i}.
 \label{eq:en-weights}
\end{equation}
The tensor product of the lattices over \(\Z_p\) is viewed as a subset of the tensor product of their fields over \(\Q_p\). Expanding the container by places gives mixed components such as \(E_1\otimes_{\Q_p}E_2\). Let \(M=\OO_{E_1}\otimes_{\Z_p}\OO_{E_2}\). The balance relation gives
\[
 (p\OO_{E_1})\otimes_{\Z_p}\OO_{E_2}=pM
 =\OO_{E_1}\otimes_{\Z_p}(p\OO_{E_2}).
\]
The index of \(p\OO_{E_i}\) in \(\OO_{E_i}\) is \(p^{d_i}\). Therefore the two presentations give \(p^{-d_1\gamma_1}\) and \(p^{-d_2\gamma_2}\). The tower law yields \(d_i\gamma_i=[k_{v_i}:\Q_p]\). For the two places above \(5\) constructed here, this gives
\[
 \Vol^\Gamma(5M)=5^{-1},\qquad \Vol^\Gamma(5M)=5^{-2}.
\]
The same contradiction applies with extra integral slots. Formula \eqref{eq:en-weights} assigns a value to a tuple of labeled factor lattices. To give a volume to the resulting tensor subset, it would have to agree on both presentations. A correction to \eqref{eq:en-printed} leaves this obstruction in place.\cite[Weighted volume and Form 2]{jiii}

For the standard extension \(|p|_{E_i}=p^{-1}\), the Haar identity is \(\mu_i(a\OO_{E_i})=|N_{E_i/\Q_p}(a)|_p=|a|_{E_i}^{d_i}\). The local degree is essential in converting between these normalizations. The argument above uses lattice indices directly.

\section{From Local Constructions to Form 2}
\subsection{Valuation scales and the period-ring supremum}
Part I uses perfectoid untilts and geometric base points to construct arithmetic holomorphic structures. Part II uses points \([a^{j^2}]-p\) of the Fargues--Fontaine curve, with \(0<|a|<1\). The valuations induced from the tilt satisfy \(v_j(p)=j^2v_1(p)\), hence \(v_j(p)=j^2v_m(p)/m^2\). An isometric embedding of the same valued field in all the factors fixes the value of \(p\) in each factor. The global normalization of Part II(1/2) restores the product formula, so an application of the local scaling must specify which valuations are used. In Part III's proof of Form 1, the exponent \(j^2/m^2\) is applied with the first factor as the reference. For \(j=1\) the displayed equality becomes \(a=a^{1/m^2}\) with \(0<a<1\) and \(m>1\).\cite{ji,jii,jhalf,jiii}

With the last factor as reference, and provided that the lifts belong to the period-ring set under consideration, the calculation is
\[
 \prod_{j=1}^m a^{j^2/m^2}=a^{(m+1)(2m+1)/(6m)}\ge a^m.
\]
Multiplication over the finite set of bad places and passage to the supremum then give the corresponding period-ring estimate. The volume in the cohomological container needs a further comparison.\cite{jiii}

\subsection{Values of the printed theta series}
After canceling the prefactor, the printed series in Part II is
\[
 \theta_M(u)=\sum_{n\in\Z}q^{n(n+1)/2}u^{2n+1}.
\]
The substitution \(n\mapsto -n-1\) gives \(\theta_M(u^{-1})=\theta_M(u)\). In odd residue characteristic the constant contribution to \(\theta_M(1)\) is \(2\), so the value is nonzero. For \(p\ne\ell\) let \(\zeta\) be a primitive \(\ell\)-th root of unity. At \(u=q^{j/(2\ell)}\zeta\) the exponent of \(q\) in term \(n\) is \(b_n=n(n+1)/2+(2n+1)j/(2\ell)\). Since \(b_{n+1}-b_n=n+1+j/\ell\), the unique minimum is at \(n=-1\) and is \(-j/(2\ell)\). The leading sum in \(\theta_M(\zeta)\) is \(\zeta+\zeta^{-1}\), a unit because the residue of \(\zeta\) has odd prime order. The ratio consequently has absolute value \(|q|^{-j/(2\ell)}\). The printed value \(|q|^{-j^2/(2\ell)}\) already differs at \(j=2\). Using a Tate root directly also requires distinguishing its absolute value \(|q|^{1/(2\ell)}\) from \(|q|\) in the same normalization.\cite{jii,jiii}

\subsection{Cohomology, tensor relations, and hulls}
For a nonzero integral Tate root \(r=q^{1/(2\ell)}\), put \(p^*=p\) when \(p\) is odd and \(p^*=4\) when \(p=2\). The compatible roots of the principal unit \(1+p^*r\) define a unit Kummer class. Its ordinary Bloch--Kato logarithm corresponds to \(\log_p(1+p^*r)\). Dividing by \(p^*\) gives a rescaled coordinate. Part III also identifies that rescaled coordinate with the inverse of the ordinary exponential. For odd \(p\) and \(u=1+p\) in \(\Q_p\), the ordinary logarithm has valuation \(1\), and the rescaled one has valuation \(0\). With these maps distinguished, the isometry used in the appendix survives.\cite{jiii}

After cohomology classes from different arithmeticoids are placed in a common space, \((u,v)\mapsto u\otimes v\) makes sense as a map of sets. It sends \((1,0)\) and \((0,1)\) to zero, while sending their sum to \(1\otimes1\), so it is not additive on the product. A cross norm with \(\|u\otimes v\|=|u|_1|v|_2\) must respect \(p\otimes1=1\otimes p\), which requires \(|p|_1=|p|_2\). This conflicts with keeping the different squared scales in the same tensor product. Cohomological isomorphisms and the global Frobenius shift can organize witnesses, whose sizes must then be computed in a common scalar normalization.\cite{jhalf,jiii,jiv,janab,jalg}

A holomorphic hull can also differ from the ordinary non-Archimedean closed convex hull. For odd \(p\), take \(E=\Q_p(\pi)\) with \(\pi^2=p\). Under \(E\otimes_{\Q_p}E\simeq E\times E\), \(a\otimes b\mapsto(ab,a\bar b)\), the integral tensor lattice has a compact, closed convex image \(P\). Its coordinates agree modulo \(\pi\) because conjugation acts trivially on the residue field. Thus \((1,0)\notin P\), although both projections of \(P\) are \(\OO_E\). Its componentwise ideal hull is \(\OO_E\times\OO_E\). In Form 2 a single witness is used to claim containment of an entire tensor lattice. Such an inference requires an additional closure operation. If \(U\) contains an integral point \(t\), its componentwise ideal hull contains \(t\OO_T\), with \(\OO_T\) the integral closure in the tensor algebra. The appendix computes a volume for this inclusion.\cite{jiii}

\subsection{The dependence of Part IV on Part III}
Part IV uses Part III's lower bound in passing from volume comparisons to heights. The printed corollary fails for the arithmetic data above, and the same data give incompatible values under Form 2's volume rule. Part IV's sign convention and shift of the standard point still require a well-defined volume and a lower bound for the actual hull. The audited ATS proof chain therefore fails to establish the claimed \(abc\) result.\cite{jiii,jiv}

\appendix
\renewcommand{\theHsection}{english.appendix.\Alph{section}}
\section{A Finite-Place Volume Lower Bound Recovered from ATS}
\subsection{Volume on the integral closure}
Let \(T\) be a tensor product over \(\Q_p\) of finitely many finite extensions of \(\Q_p\). As a finite étale algebra it is a product \(\prod_\alpha F_\alpha\), with integral closure \(\OO_T=\prod_\alpha\OO_{F_\alpha}\). Normalize additive Haar measure by \(\mu_T(\OO_T)=1\) and put \(\nu_T(U)=(\dim_{\Q_p}T)^{-1}\log\mu_T(U)\). For an integral set \(U\) with a nonzero projection in each component, its smallest componentwise ideal hull is
\[
 H(U)=\prod_\alpha\pi_\alpha^{n_\alpha}\OO_{F_\alpha},
 \qquad n_\alpha=\min_{u\in U}\ord_\alpha(u_\alpha).
\]
The nonzero valuations are discrete and bounded below, so the minima exist. For a singleton \(t\) the hull is \(t\OO_T\). If \(T=\bigotimes_i E_i\) and \(t=\bigotimes_i t_i\) with each \(t_i\ne0\), the determinant of multiplication gives
\begin{equation}
 \nu_T(t\OO_T)=\sum_i\log|t_i|_{\rm st},\qquad |p|_{\rm st}=p^{-1}.
 \label{eq:en-haar}
\end{equation}
Indeed, the determinant from slot \(i\) is \(N_{E_i/\Q_p}(t_i)^{\dim T/[E_i:\Q_p]}\). The normalized determinant is intrinsic to the multiplication map on \(T\).

\subsection{A Kummer witness in one active slot}
Let \(K/F\) be a finite extension of number fields whose completions contain the required local Tate roots, and let \(B_F\) be the chosen finite set of bad places. For \(s\mid w\in B_F\) take \(r_s^{2\ell}=q_w\), and put \(r_s=1\) elsewhere. Define
\[
 \tau_p(r)=p^{*-1}\log_p(1+p^*r).
\]
For nonzero integral \(r\), we have \(v_p(p^*r)>1/(p-1)\). Since \(v_p(n)\le(n-1)/(p-1)\), the linear term of the logarithm strictly dominates all the higher terms. Hence \(|\tau_p(r)|_{\rm st}=|r|_{\rm st}\), and \(\tau_p(1)\) is a unit.

\begin{trivlist}\item[\hskip\labelsep\textbf{Proposition 1.}]
In label group \(j\), use \(j+1\) slots, with \(\tau_p(r_s)\) in slot \(j\) and \(\tau_p(1)\) in the others. Take the componentwise principal ideal hulls of these points and average the dimension-normalized log-volumes over the label groups. The finite-place average is \(-Q\), where
\[
 Q=\sum_p Q_p,\qquad
 Q_p=\frac{\log p}{2\ell[F:\Q]}
       \sum_{\substack{w\mid p\\w\in B_F}} f_w\ord_w(q_w)
\]
and \(f_w\) is the residue degree of \(F_w/\Q_p\). Any finite-place region containing these hulls has the corresponding lower bound when its global volume is defined.
\end{trivlist}
\begin{proof}
Write \(A_p=K\otimes_{\Q}\Q_p=\prod_{s\mid p}K_s\) and \(\lambda_s=[K_s:\Q_p]/[K:\Q]\). Expand \(A_p^{\otimes(j+1)}\) by place tuples \(c=(s_0,\ldots,s_j)\). On component \(c\), equation \eqref{eq:en-haar} gives normalized log-volume \(\log|r_{s_j}|_{\rm st}\). The dimension proportion of this component is \(\prod_i\lambda_{s_i}\). Summing over the filler slots uses \(\sum_{s\mid p}\lambda_s=1\), leaving
\[
 h_{j,p}=\sum_{s\mid p}\lambda_s\log|r_s|_{\rm st}.
\]
For \(w\in B_F\) let \(e_w\) be the ramification index of \(F_w/\Q_p\). Then
\[
 -\log|r_s|_{\rm st}=\frac{\ord_w(q_w)}{2\ell e_w}\log p,\qquad
 \sum_{s\mid w}\lambda_s=\frac{e_w f_w}{[F:\Q]}.
\]
The second identity follows from \(\sum_{s\mid w}[K_s:F_w]=[K:F]\). Thus \(h_{j,p}=-Q_p\) for each label group. Summation over \(p\) and averaging over the labels give \(-Q\). All other places contribute zero. The lower bound for containing regions follows by monotonicity.
\end{proof}

\subsection{Comparison with the concrete reference region}
The concrete reference region uses \(r_s^{j^2}\) in the active slot. Its allowed local field automorphisms preserve the associated principal ideal. For \(j>1\) and \(0<|r_s|_{\rm st}<1\), one has \(r_s^{j^2}\OO\subsetneq r_s\OO\), and the Kummer witness lies outside the smaller ideal in that component. Dimension weighting gives the reference finite-place value \(-j^2Q_p\). With the concrete archimedean contribution \((j+1)\log\pi\), the average global value is
\[
 R_{\rm ref}=-A_m Q+B_m,\qquad
 A_m=\frac{(m+1)(2m+1)}6,\qquad B_m=\frac{m+3}2\log\pi .
\]
Thus \(-Q\le R_{\rm ref}\) is equivalent to \((A_m-1)Q\le B_m\). For \(\ell=2m+1\ge7\) this requires \(2\ell Q\le(126/11)\log\pi\). For the rational Legendre parameter \(x_N=5^N/(2\cdot5^N+1)\), the \(j\)-invariant has a pole of order \(2N\) at \(5\). If all places above \(5\) belong to the support, then \(2\ell Q\ge2N\log5\), which exceeds the threshold for \(N\ge6\). The local Lean existence theorem for this family takes the curve data, a finiteness hypothesis, and an anabelian existence hypothesis as explicit inputs.\cite{localref,adapt}

The ATS Ansatz permits trying a separate choice of standard point for the active slot of each label group. To use these witnesses, one must place them in the very hull measured by the upper bound and account for the archimedean terms and the place weights. For descent to a smaller local field \(k\), the trace of the entire ideal can be computed. If \(\mathfrak D_{E/k}=\mathfrak p_E^d\), then
\[
 {\rm Tr}_{E/k}(\mathfrak p_E^n)
 =\mathfrak p_k^{\lfloor(n+d)/e(E/k)\rfloor}.
\]
Indeed, the trace lands in \(\mathfrak p_k^a\) exactly when \(\mathfrak p_E^{n-ea}\subseteq\mathfrak D_{E/k}^{-1}\), by trace duality. For \(b\OO_E=\mathfrak p_E^d\), applying the trace to \(b^{-1}\mathfrak p_E^n\) and restricting its image to \(\mathfrak p_k^{\lceil n/e\rceil}\) adds less than \(1\) to the exponent \(n/e\). Weighted over the finite bad places, this rounding loss is bounded by a log radical term. The calculation applies to the entire ideal. A single trace can lose valuation information through cancellation, and the original tuple-wise upper estimate still needs a new justification for the revised ideal.\cite{riemcode,adapt}

\subsection{What is needed for Szpiro and \texorpdfstring{\(abc\)}{abc}}
The new lower bound retains the multiplicities of the Tate parameters. Fix an odd prime \(p\) and take \(a=p^N\), \(b=1\), and \(c=p^N+1\). The Frey curve \(y^2=x(x-a)(x+b)\) has minimal discriminant valuation and Tate parameter valuation \(2N\) at \(p\). The contribution of \(p\) to \(\log\rad(abc)\) is always \(\log p\). The prime factors of \(p^N+1\) also enter the global radical. A Szpiro-type bound for a family still requires an estimate relating such multiplicities to the conductor or radical.

Suppose estimates on the same actual hull give
\[
 -Q-\delta\le R\le-AQ+U,\qquad A>1,\quad\delta\ge0.
\]
Then \((A-1)Q\le U+\delta\). Converting the right-hand side to radical and height terms uniformly on the applicable family can yield a Szpiro or \(abc\) estimate on that family. A fixed exponent estimate for all coprime positive triples \(a+b=c\) would have the form \(\log c\le B\log\rad(abc)+C\) with fixed \(B>1\) and \(C\). Full \(abc\) requires \(B=1+\varepsilon\) for every \(\varepsilon>0\), with \(C_\varepsilon\) uniform over all triples. The finite-place lower bound proved here leaves these comparisons and uniformity requirements open.

For the reference coefficient \(A=A_m=\ell(\ell+1)/12\), put \(L=2\ell Q\). Eliminating \(R\) yields
\[
 \frac L6\le\frac{4\ell}{(\ell-3)(\ell+4)}(U+\delta).
\]
For rational Frey curves with the chosen bad places above \(2abc\), the whole-ideal rounding error is at most \(\log\rad(abc)+\log 2\). Its coefficient in this estimate is \(O(1/\ell)\), so this error allows the exponent to approach \(1\) if sufficiently large admissible auxiliary primes can be chosen uniformly. The lower and upper estimates for the actual hull and the arithmetic bound for \(U\) still determine the attainable exponent.

\section[Theta Lifts and ATS]{Using and Extending the Theta-Lift Techniques behind Dirichlet GRH to Further Salvage ATS}

On 21 November 2026, the Lean community blog announced that it had verified Levent Alpöge's proof of the classical Riemann hypothesis. The Riemann project's Lean source also gives a theorem for Dirichlet \(L\)-functions stating that their zeros in the critical strip have real part \(1/2\).\cite{riemblog,riemcode}

The project's trace duality and different formulas apply to ideal descent in Appendix A. Its explicit prime counts help control the error in choosing the auxiliary prime \(\ell\). Both applications are independent of the zero-location theorems. Here we examine whether the project's automorphic theta lifts and Fourier techniques can also bound local volume deficits in ATS. Applying such a bound to Szpiro or \(abc\) still requires realizing the Kummer witnesses in the actual hull and establishing a global upper bound on that same hull.\cite{riemcode,riemtheta}

\subsection{Local witnesses from consecutive theta values}
At an odd bad place \(p\ne\ell\), work in one local field \(E\) containing \(r=q^{1/(2\ell)}\) and a primitive \(\ell\)-th root of unity \(\zeta\). Using the series printed in Joshi II, define
\[
 \xi_0=1,\qquad \xi_j=\frac{\theta_M(\zeta)}{\theta_M(r^j\zeta)},\qquad
 u_j=\xi_j/\xi_{j-1}.
\]
The leading-term calculation gives \(v(\xi_j)=jv(r)\) and \(v(u_j)=v(r)\). The logarithm isometry of Appendix A gives \(\tau_p(u_j)\OO_E=r\OO_E\). This supplies a generator from theta values. Applying it to ATS requires the Kummer class of \(1+p^*u_j\) to reach the specified label under the gathering maps allowed in Part III. Subtraction of two Kummer classes corresponds to taking the quotient of their principal units. A quotient across untilts also requires a compatible common field and markings.\cite{jii,jiii}

\subsection{An explicit local generation calculation}
Take \(p=29,\ell=7,q=29,\pi^{14}=29\), and a primitive seventh root \(\zeta\in\Q_{29}\). In \(E=\Q_{29}(\pi)\), put \(W=\pi\OO_E\) and \(L_0=\pi^4\OO_E\) for label \(j=2\). The quotient \(W/L_0\) has dimension three over \(\mathbb F_{29}\). The rescaled logarithms of all fourteen branches \(\omega\pi\) lie on the line \(\mathbb F_{29}\pi\) modulo \(L_0\), since \(\mu_{14}\subset\Q_{29}\). Together with \(L_0\), they generate a lattice of index \(29^2\) in \(W\).

Expanding the series modulo \(\pi^4\) gives
\[
 \tau_{29}(\pi)\equiv\pi,\qquad
 \tau_{29}(\xi_2)\equiv(1+\zeta^2)\pi^2,\qquad
 \tau_{29}(u_2)-\tau_{29}(\pi)\equiv\zeta^2\pi^3.
\]
The last congruence follows from \(\theta_M(\pi\zeta)/\theta_M(\pi^2\zeta)\equiv\pi+\zeta^2\pi^3\). Higher terms of the rescaled logarithm vanish modulo \(\pi^4\), and \(1+\zeta^2\) is a unit. The three specified points and \(L_0\) therefore generate \(W\). The application requires realizing these points in the same actual label and controlling the other tensor slots. For \(r_t=29^t\pi\) with \(t\ge1\), the inclusion \(r_t^4\OO_E\subset29r_t\OO_E\) leaves fourteen independent directions in \(r_t\OO_E/29r_t\OO_E\). Three points cannot generate that quotient.

\subsection{Measuring the volume deficit by Fourier duality}
In the full \((j,p)\) tensor algebra, let \(d\) be the dimension and use the volume normalization of Appendix A. The first-power lattice \(W\) and squared-weight lattice \(L_0\) have \(\nu(W)=-Q_p\) and \(\nu(L_0)=-j^2Q_p\). Suppose the actual local region \(C\) is a measurable additive subgroup containing \(L_0\). Set \(H=(C\cap W)/L_0\subset W/L_0\). Finite group duality gives \(|H^\perp|=[W:C\cap W]\), so
\[
 \nu(C)\ge -Q_p-\frac{\log|H^\perp|}{d}.
\]
For a general measurable region with positive density \(\rho=\mu(C\cap W)/\mu(W)\), one likewise has \(\nu(C)\ge -Q_p+d^{-1}\log\rho\). Either estimate permits summing the actual local losses in the global normalization.

The Riemann project provides Heisenberg translation expansions, the Fourier formula for the Weil representation, and trace duality for working with prescribed local cosets and characters. Under its local hypotheses, the existing Rallis formula selects a test polynomial, measure, fundamental domain, and constant existentially and computes the squared Petersson norm of the lift. To use it here, one needs prescribed test data constructed from the Kummer images allowed by Part III and a theorem relating the Fourier coefficients to \(|H^\perp|\) or \(\rho\). Polarization to extract a mixed inner product would also require controlling the error after subtracting two norms in the same normalization. The zero-location conclusion of Dirichlet GRH supplies no such comparison.\cite{riemtheta}

\subsection{Linear weights preserve the leading coefficient}
If only \(\xi_j\) can be realized, then \(\tau_p(\xi_j)\OO_E=r^j\OO_E\). Degree weighting in the one-active-slot model gives \(-jQ_p\) when the witnesses exist at each selected bad place. Put \(b_m=(\ell+1)/4\) and \(A_m=\ell(\ell+1)/12\) for the averages of the linear and squared weights. If both bounds hold on the same actual global hull, then for \(L=2\ell Q\) and \(\ell>3\),
\[
 -b_mQ-\delta\le R\le-A_mQ+U
 \quad\Longrightarrow\quad
 \frac L6\le\frac{4\ell}{(\ell-3)(\ell+1)}(U+\delta).
\]
The coefficient increases by a factor of \((\ell+4)/(\ell+1)\) relative to the first-power bound in Appendix A, preserving the leading term \(4/\ell\).

For the specified tuple at \(p=337,\ell=7,j=2\), the linear lattice has log-volume \(-(6N+30/7)\log337\), and the existing tuple-wise upper bound is \(-(12N+5)\log337\). The deficit \((6N+5/7)\log337\) grows without bound. A Szpiro-type estimate by this route still requires control of which witnesses the actual hull contains and a global upper bound for that hull or compensating contributions at other places.

\section*{Use of Artificial Intelligence and License}
Edison Mai led the research and is the corresponding author. OpenAI ChatGPT (Astra 6.2) assisted with reading the papers and code, developing and organizing the counterexamples and the revised natural-language argument, and drafting, typesetting, and checking the paper. This paper and the natural-language materials are licensed under \href{https://creativecommons.org/licenses/by-nc-sa/4.0/}{CC BY-NC-SA 4.0}. The companion code is licensed under the \href{https://www.apache.org/licenses/LICENSE-2.0}{Apache License 2.0}.
\endgroup

\begingroup
\zihao{5}\fontsize{10.5bp}{16bp}\selectfont
\renewcommand{\refname}{参考文献}
\setlength{\parskip}{0bp}
\begin{thebibliography}{99}
\setlength{\itemsep}{3bp}
\setlength{\parsep}{0bp}
\bibitem{mi} Shinichi Mochizuki. Inter-universal Teichmüller Theory I. Construction of Hodge Theaters. \emph{Publ. Res. Inst. Math. Sci.} 57 (2021). \url{https://doi.org/10.4171/prims/57-1-1}
\bibitem{mii} Shinichi Mochizuki. Inter-universal Teichmüller Theory II. Hodge--Arakelov-theoretic Evaluation. \emph{Publ. Res. Inst. Math. Sci.} 57 (2021). \url{https://doi.org/10.4171/prims/57-1-2}
\bibitem{miii} Shinichi Mochizuki. Inter-universal Teichmüller Theory III. Canonical Splittings of the Log-theta-lattice. \emph{Publ. Res. Inst. Math. Sci.} 57 (2021). \url{https://doi.org/10.4171/prims/57-1-3}
\bibitem{miv} Shinichi Mochizuki. Inter-universal Teichmüller Theory IV. Log-volume Computations and Set-theoretic Foundations. \emph{Publ. Res. Inst. Math. Sci.} 57 (2021). \url{https://doi.org/10.4171/prims/57-1-4}
\bibitem{ji} Kirti Joshi. Constructions of Arithmetic Teichmuller Spaces I. \url{https://arxiv.org/abs/2106.11452}
\bibitem{jii} Kirti Joshi. Constructions of Arithmetic Teichmuller Spaces II. \url{https://arxiv.org/abs/2303.01662}
\bibitem{jhalf} Kirti Joshi. Constructions of Arithmetic Teichmuller Spaces II(1/2). \url{https://arxiv.org/abs/2305.10398}
\bibitem{jiii} Kirti Joshi. Constructions of Arithmetic Teichmuller Spaces III. \url{https://arxiv.org/abs/2401.13508v4}. arXiv v4 的 TeX 源文件与仓库所收文件一致，SHA-256 前 16 位为 \nolinkurl{b3df2cfb3107c171}。可复核式 \eqref{eq:printed} 于源文件第 3649 行，带权定义及实例于第 3529、3554 行。
\bibitem{jiv} Kirti Joshi. Constructions of Arithmetic Teichmuller Spaces IV. \url{https://arxiv.org/abs/2403.10430}
\bibitem{janab} Kirti Joshi. On Mochizuki's Idea of Anabelomorphy and Its Applications. \url{https://arxiv.org/abs/2003.01890}
\bibitem{jalg} Kirti Joshi. Algebraization of Mochizuki's Anabelian Variation of Ring Structures, Perfectoid Geometry and Formal Groups. \url{https://arxiv.org/abs/1906.06840}
\bibitem{serre} Jean-Pierre Serre. Propriétés galoisiennes des points d'ordre fini des courbes elliptiques. \emph{Invent. Math.} 15 (1972), 259--331. \url{https://doi.org/10.1007/BF01405086}
\bibitem{lana} LANA Project. IUT IV Section 1 Formalization. \url{https://github.com/LANA-Project/iut4-sec1}. 相对路径 \lana{README.md}{README.md}、\lana{Plans/Iut4Sec1Spec.md}{Plans/Iut4Sec1Spec.md}、\lana{blueprint-verso/Iut4Sec1Blueprint/Chapters/Theorem110.lean}{blueprint-verso/Iut4Sec1Blueprint/Chapters/Theorem110.lean}。
\bibitem{localref} 本项目使用的本地参考实现，相对本项目根目录 \nolinkurl{refer/iut/Iut/Tripod/ClassicalAbcGenuineCanLift.lean:40}、\nolinkurl{refer/iut/Iut/Cor312/Statement.lean:39}、\nolinkurl{refer/iut/Iut/Concrete/ThetaRegion.lean:174}。项目仓库 \url{https://github.com/NanachiKyubey/joshi-teich-abc}。
\bibitem{adapt} 本项目的 Lean 适配代码。相对路径 \project{prove2me_workspace/Solutions/Source_ReferenceIutCor312GlobalVolume.lean\#L47}{prove2me_workspace/Solutions/Source_ReferenceIutCor312GlobalVolume.lean:47}、\project{prove2me_workspace/Solutions/Source_TripodRationalLegendreFiveP7Counterexample.lean\#L197}{prove2me_workspace/Solutions/Source_TripodRationalLegendreFiveP7Counterexample.lean:197}。
\bibitem{riemblog} Lean Community. Lean Verification of the Riemann Hypothesis. 21 November 2026. \url{https://leanprover-community.github.io/blog/posts/riemann-hypothesis/}
\bibitem{riemcode} Levent Alpöge. Riemann Lean 项目压缩包。以下路径相对于解压后的 Riemann 项目根目录。\par
\nolinkurl{src/Main.lean:4}\par
\nolinkurl{src/Statements.lean:3}\par
\nolinkurl{src/GroundTheory/LocalFields/LocalFourierTrace.lean:28}\par
\nolinkurl{src/GroundTheory/LocalFields/Different.lean:132}\par
\nolinkurl{src/Formalism/PeriodIntegrals/OrbitalBoundsConclusion/PrimeCount.lean:7}\par
\nolinkurl{src/Formalism/PeriodIntegrals/PrimesDyadicRS.lean:5}
\bibitem{riemtheta} Levent Alpöge. Riemann Lean archive, theta lifting and Fourier analysis. Paths relative to the archive root.\par
\nolinkurl{src/Development/AbelianVarietyTheory/Proofs/DecoratedLatticeProofs/ThetaFourierSumEqSumHeisenbergTranslates.lean:20}\par
\nolinkurl{src/Development/AbelianVarietyTheory/Proofs/DecoratedLatticeProofs/WeylElementActsOnSchrodingerByFourier.lean:40}\par
\nolinkurl{src/Development/AutomorphicArithmetic/ThetaCorrespondence/ThetaPairLifts/LocalTypeDictionary/BigThetaNormEqRallisInnerProduct.lean:19}\par
\nolinkurl{src/GroundTheory/LocalFields/LocalFourierTrace.lean:28}
\end{thebibliography}
\endgroup
\end{document}
